\subsection{Baseline Fidelity and Evaluation Scope}

The primary CaT and GECC rows retain feature distribution matching and
clustering of propagated features, respectively, under the common
continual access protocol. Both attach text from five source nodes and use
the common student. This text attachment is an adaptation of the original
numerical algorithms. GCond-x uses the official simplified graph convolution
(SGC) model and gradient distance objective with identity adjacency and a
changed continual proxy. SFGC uses the official trajectory matching
objective with reduced teacher and construction budgets.
The final BookHis increment has one proxy class, giving no discriminative
feature gradient for these two classification-based cores.

The complete five-dataset comparison, including the selection controls and the
two small-graph results, is consolidated in \autoref{tab:main_comparison}.
The appendix therefore retains the protocol and fidelity notes rather than a
second copy of the same numerical table.

TAGSAM retains the official selector/response core under a bounded budget
of 100 outer iterations and 10 response minibatch updates per iteration.
The official configuration uses 5,000 outer iterations with 15 full
synthetic-data epochs per iteration. An epoch traverses the synthetic
dataset, so these response budgets are not interchangeable update counts.
STAD uses a local Qwen generator with changed supervision, synthetic
topology, and continual access. These adapted baselines in
\autoref{tab:matched_main} must not be read as complete native
reproductions. GECC's native soft fuzzy c-means (FCM) core completed BookHis
only (40.54 F1, 74.26 Accuracy).
Its other large-graph results are unavailable.
Earlier controls inspired by DANCE and SR-GM, and frozen GraphCLIP checkpoints,
use different students or task definitions and do not enter the common student
ranking. Cora/Citeseer use the registered project split, not the Planetoid
protocol with 20 labeled training nodes per class.

For Cora and Citeseer, TAGSAM keeps GPT-2-large text selection based on
pointwise mutual information (PMI) and the reduced response training budget. STAD keeps native coarsening,
self-training, and internal candidate validation, with local
Qwen2.5-1.5B generation and binary stored support for the common student model.
Its internal split assigns approximately 10\% of each current class to
training and 10\% to validation, reserving the remainder for pseudo-labeling.
external evaluation labels do not select text candidates.
The earlier inspired controls use a contrastive student with one positive
text per node, 1,024 updates, a three-stage ArXiv stream, and the development
partition. The common comparison instead treats texts from the same class
as positives, adds class prompt supervision, trains for 2,048 updates, and
evaluates the historical test partition. \autoref{tab:earlier_references}
preserves the earlier scores separately because both training and evaluation
differ. Comparing these controls under the current protocol would require
recovering or reconstructing their memories under the declared stage
schedule and training the common student. They would remain inspired
implementations rather than official reproductions.
SR-GM-inspired projects conflicting
gradients in a common 128-dimensional latent space because the original
feature and text dimensions differ. It is not an official SR-GM reproduction.
Completed repetitions in \autoref{tab:main_comparison} cross three class
orders with three construction seeds, each averaging three student runs.
Every main-table cell therefore reports nine memory means. Its SD is the
sample standard deviation across those nine memory means.

\begin{table*}[t]
\centering
\scriptsize
\caption{Earlier inspired references, kept separate from the common
comparison. Entries are Macro-F1 / Accuracy on the development partition, mean $\pm$ SD
over three student model seeds for one order and construction seed. These use
1,024 updates, one positive text per node in the contrastive loss, and a three-stage
ArXiv stream. These implementations differ from the original methods.}
\label{tab:earlier_references}
\resizebox{\textwidth}{!}{%
\begin{tabular}{@{}lccccc@{}}
\toprule
Method & ArXiv & SportsFit & BookHis & Cora & Citeseer\\
\midrule
DANCE-inspired & 17.61$\pm$0.72 / 23.59$\pm$2.27
& 33.07$\pm$2.85 / 40.11$\pm$4.75 & 14.68$\pm$0.56 / 20.51$\pm$1.11
& 39.83$\pm$0.55 / 47.11$\pm$1.54 & 43.63$\pm$1.42 / 49.84$\pm$1.44\\
SR-GM-inspired & 20.97$\pm$0.58 / 32.13$\pm$2.12
& 30.89$\pm$4.87 / 37.22$\pm$11.39 & 15.74$\pm$1.87 / 26.43$\pm$4.79
& 28.41$\pm$4.38 / 31.12$\pm$6.07 & 48.94$\pm$3.02 / 53.50$\pm$3.63\\
STAD-inspired & 17.69$\pm$0.78 / 24.70$\pm$2.18
& 28.71$\pm$3.74 / 33.77$\pm$4.98 & 15.91$\pm$0.57 / 22.88$\pm$2.31
& 29.84$\pm$1.79 / 33.95$\pm$2.88 & 46.08$\pm$1.66 / 52.56$\pm$1.99\\
\bottomrule
\end{tabular}}
\end{table*}

\subsection{ArXiv Protocol and Additional Experiments}
\label{app:arxiv_archive}

The main ArXiv class arrival protocol uses three stages with 13, 14, and
13 classes. The stage partition determines the current graph, its
neighborhoods, and each cumulative memory budget, so all main ArXiv rows use
this schedule. \autoref{tab:arxiv_three_stage} repeats the ArXiv column for
auditability.

\begin{table}[H]
\centering
\footnotesize
\caption{ArXiv comparison repeated for auditability. The three-stage class schedule is 13/14/13, matching the ArXiv column in \autoref{tab:main_comparison}. Entries are mean $\pm$ SD across nine memories. Each memory is evaluated with three student seeds.}
\label{tab:arxiv_three_stage}
\begin{tabular*}{\linewidth}{@{\extracolsep{\fill}}lcr@{}}
\toprule
Method & $n$ & Macro-F1 / Accuracy\\
\midrule
CaT & 9 & 34.98$\pm$1.79 / 60.41$\pm$2.64\\
GECC & 9 & 52.68$\pm$0.45 / 70.94$\pm$0.49\\
Random & 9 & 47.72$\pm$0.56 / 69.39$\pm$0.30\\
K-center & 9 & 30.84$\pm$1.50 / 59.24$\pm$1.26\\
Cluster medoid & 9 & 48.30$\pm$0.86 / 69.41$\pm$0.27\\
Herding & 9 & 49.21$\pm$0.18 / 69.63$\pm$0.34\\
GCond-x & 9 & 49.64$\pm$1.28 / 66.71$\pm$0.93\\
SFGC & 9 & 24.23$\pm$2.17 / 49.10$\pm$3.41\\
TAGSAM & 9 & 39.31$\pm$3.11 / 58.92$\pm$3.54\\
STAD & 9 & 49.76$\pm$0.64 / 66.07$\pm$1.37\\
\method{} & 9 & 52.91$\pm$0.73 / 71.15$\pm$0.59\\
\bottomrule
\end{tabular*}
\end{table}


An earlier five-stage study evaluated only the development partition
(\autoref{tab:arxiv_confirmation}). These scores are separate development
evidence, not the main ArXiv test comparison, and its memories were not
retained for reevaluation.

\begin{table}[H]
\centering
\footnotesize
\setlength{\tabcolsep}{3pt}
\caption{Earlier ArXiv five-stage study on the development partition.
Values are mean $\pm$ sample SD across nine memories from three class
orders and three construction seeds, each averaging three student runs.
These are not the test scores for \autoref{tab:main_comparison}.}
\label{tab:arxiv_confirmation}
\begin{tabular*}{\linewidth}{@{\extracolsep{\fill}}lrr@{}}
\toprule
Method & Macro-F1 & Accuracy\\
\midrule
CaT & 35.32$\pm$2.11 & 62.75$\pm$2.21\\
GECC & 51.31$\pm$0.74 & 71.10$\pm$0.57\\
\method{} & \textbf{51.93$\pm$1.06} & \textbf{71.65$\pm$0.40}\\
\bottomrule
\end{tabular*}
\end{table}

The paired Macro-F1 difference is +0.62 points against GECC
(95\% interval $[-0.45,1.69]$, positive in six of nine combinations)
and +16.61 against CaT ($[15.46,17.77]$, positive in all nine).
This supports utility under a longer stream without establishing a
reliable F1 advantage over GECC on ArXiv.

\autoref{tab:arxiv_archived_controls} retains the earlier ArXiv mechanism
results, all using the three-stage schedule. The feature and assignment
comparisons are separate batches and are not interchangeable.

\begin{table}[H]
\centering
\footnotesize
\caption{Archived ArXiv controls with three stages. Each construction batch uses one class order and three construction seeds, each averaging three student runs.}
\label{tab:arxiv_archived_controls}
\begin{tabular*}{\linewidth}{@{\extracolsep{\fill}}lrr@{}}
\toprule
Loss control & F1 & Accuracy\\
\midrule
\method{} & 52.96 & 70.64\\
Fixed $Q$ & 40.20 & 55.80\\
Without DM & 51.05 & 64.19\\
Without GM & 53.20 & 70.83\\
\midrule
Feature construction & GCN F1 & GraphSAGE F1\\
\midrule
Aggregation through $Q$ & 52.96$\pm$0.21 & 54.02$\pm$0.69\\
Direct optimization & 52.27$\pm$0.48 & 51.74$\pm$0.49\\
\midrule
Assignment & F1 & GiB / seconds\\
\midrule
Separate & 53.20$\pm$0.12 & 15.99 / 292.5\\
Shared & 53.24$\pm$0.09 & 12.08 / 241.3\\
\bottomrule
\end{tabular*}
\end{table}

The earlier ArXiv full-history reference gives 53.47 F1, compared with
52.96 for the memory, corresponding to 99.04\% F1 retention and 97.60\%
Accuracy retention. Student training alone takes 100.7 MiB and 7.75 s with
memory, compared with 1,614.6 MiB and 63.61 s with all construction nodes.
These use the same measurement boundaries as
\autoref{tab:retention} and~\autoref{tab:student_resource_audit}, with the
earlier ArXiv stage schedule.

\autoref{fig:split_check} likewise belongs to the earlier comparison.
Its historical labels refer to the CaT and GECC adaptations described above.
The appendix tables below use the same three-stage ArXiv class schedule when
they report class-incremental controls. Chronological results and the separate
five-stage development check are identified explicitly.

\begin{figure}[t]
\centering
\includegraphics[width=\linewidth]{primary_confirmation.pdf}
\caption{Archived relative Macro-F1 differences from GECC, with 95\% confidence intervals over nine matched memories. ArXiv uses three stages. Other datasets use the schedules in \autoref{tab:datasets}.}
\Description{Relative Macro-F1 differences between CoTAG-Q and GECC for five datasets. ArXiv uses the three-stage class schedule.}
\label{fig:split_check}
\end{figure}

\subsection{Chronological Protocol and Prediction Scope}
\label{app:temporal_setting}

The chronological ArXiv stream uses publication-year cutoffs of 2015,
2016, 2017, 2018, and 2020. ACM and DBLP are constructed citation subsets
of the AMiner ACM v9 and DBLP v10 archives. ACM contains 16,928 nodes in
three venue-derived fields (databases, data mining, and networking), with
11,848 construction nodes and 592 final memory items. DBLP contains
57,520 nodes in four fields (artificial intelligence (AI), computer vision,
natural language processing (NLP), and machine learning), with 40,227
construction nodes and 2,011 memory items.
The training-year quantiles at 40\%, 55\%, 70\%, 85\%, and 100\%
give cutoffs 2005/2008/2010/2012/2015 for ACM and
2008/2010/2012/2014/2015 for DBLP. Papers from the same year stay
together. Five exact normalized-title overlaps with ACM were removed
from DBLP. These are curated subsets rather than unchanged official
benchmark splits.

All methods construct each new block from the current arrival group
and its induced edges, retain the historical synthetic prefix, and use
the cumulative 5\% budget. The complete 30-update \method{} objective
is retained. Its chronological runs use float64 Sinkhorn normalization
with an iteration cap of 16,384 instead of 4,096, tolerance $10^{-4}$,
and float32 output. This recorded numerical repair resolves a failed
transport solve without adding outer construction updates.
The year order is fixed, with three construction seeds and three
student runs per memory. The central processing unit (CPU) loader retains the source graph, so
the access protocol is enforced within the constructor rather than by
a streaming loader.

\autoref{tab:temporal_results} uses all observed test nodes and the
cumulative graph at the final stage. The paired F1 differences from
GECC are +0.75 points on ArXiv (95\% interval [0.10, 1.40]),
-0.51 on ACM ([-1.75, 0.74]), and -0.74 on DBLP ([-3.54, 2.06]).
\autoref{tab:temporal_current} instead restricts both queries and
inference edges to the current arrival group. It assesses prediction
when historical graph context is unavailable. Because the query
cohorts also differ, the difference between these tables does not
isolate the effect of historical edges.

\begin{table}[H]
\centering
\footnotesize
\setlength{\tabcolsep}{3pt}
\caption{Final-stage Macro-F1 with queries and inference graph restricted
to the current arrival group. Values average three construction seeds,
each averaging three student runs.}
\label{tab:temporal_current}
\begin{tabular*}{\linewidth}{@{\extracolsep{\fill}}lrrr@{}}
\toprule
Method & ArXiv & ACM & DBLP\\
\midrule
CaT & 37.46 & 85.14 & 76.59\\
GECC & 48.10 & 85.10 & 75.30\\
\method{} & 49.83 & 85.18 & 73.75\\
\bottomrule
\end{tabular*}
\end{table}

On DBLP, mean F1 over the five stages when predicting only current arrivals
differs by -0.91 points from GECC (95\% interval [-1.57, -0.26]) and
-2.76 from CaT ([-4.74, -0.78]). The chronological evidence therefore supports
operation with recurring classes, with predictive gains depending
on the dataset and available graph context.

\subsection{Additional Evaluation Partitions and Continual Metrics}

\begin{table}[H]
\centering
\footnotesize
\setlength{\tabcolsep}{3pt}
\caption{Additional evaluation partition, with one order, three construction seeds, and
three student model seeds per constructed memory. Each entry is mean \(\pm\)
sample SD over the three construction means.}
\label{tab:additional_mask}
\begin{tabular*}{\linewidth}{@{\extracolsep{\fill}}llcc@{}}
\toprule
Dataset & Method & Macro-F1 & Accuracy\\
\midrule
ArXiv & \method{} & 53.40$\pm$0.13 & 71.64$\pm$0.25\\
& GECC & 51.89$\pm$0.63 & 71.28$\pm$0.67\\
& CaT & 35.70$\pm$0.62 & 62.15$\pm$1.13\\
\midrule
SportsFit & \method{} & 79.56$\pm$0.24 & 87.45$\pm$0.60\\
& GECC & 80.30$\pm$0.55 & 89.25$\pm$0.12\\
& CaT & 75.04$\pm$1.90 & 86.14$\pm$1.96\\
\midrule
BookHis & \method{} & 45.16$\pm$2.03 & 81.33$\pm$0.32\\
& GECC & 42.01$\pm$0.24 & 82.21$\pm$0.25\\
& CaT & 26.87$\pm$0.99 & 80.17$\pm$0.39\\
\bottomrule
\end{tabular*}
\end{table}

Against GECC, F1 differences on the additional partition have paired intervals
of [-0.24, 3.25] on ArXiv, [-2.68, 1.19] on SportsFit, and
[-2.47, 8.77] on BookHis. The respective positive paired-difference counts are
3/3, 1/3, and 3/3.
This experiment uses three construction seeds and an altered construction
pool, whereas the main comparison crosses three orders and three seeds.

\begin{table}[H]
\centering
\footnotesize
\setlength{\tabcolsep}{3pt}
\caption{Continual diagnostics over nine combinations of class order and construction seed.
Old and New report final-stage Macro-F1 for classes introduced before and
during the final stage, respectively. Stage reports mean Macro-F1 across
stages. ArXiv uses the three-stage schedule in \autoref{tab:datasets}. The
other datasets match \autoref{tab:matched_main}.}
\label{tab:continual_metrics}
\begin{tabular*}{\linewidth}{@{\extracolsep{\fill}}llrrr@{}}
\toprule
Dataset & Method & Old & New & Stage\\
\midrule
ArXiv & \method{} & 52.25 & 54.27 & 62.23\\
& GECC & 51.45 & 55.24 & 62.34\\
& CaT & 33.77 & 37.50 & 47.36\\
\midrule
SportsFit & \method{} & 80.00 & 80.51 & 82.85\\
& GECC & 80.57 & 80.58 & 83.59\\
& CaT & 74.31 & 75.52 & 79.50\\
\midrule
BookHis & \method{} & 48.39 & 53.45 & 51.82\\
& GECC & 45.48 & 43.86 & 50.43\\
& CaT & 29.16 & 29.09 & 34.87\\
\bottomrule
\end{tabular*}
\end{table}

On BookHis, the gain over GECC is larger for new classes than for old
classes. On ArXiv and SportsFit, the stage averages remain close to GECC.

\subsection{Architecture Transfer}

\autoref{tab:sage} reports both student models for the three-stage ArXiv
schedule. The constructor is not rerun and memory hashes match across
architectures. GraphSAGE uses two mean SAGEConv layers with explicit
self-loop neighbors and the root transform. Against GECC, \method{} gains
1.34 F1 points with positive differences in all nine combinations. An exploratory
bootstrap that resamples orders and construction seeds separately gives an
interval of [0.74, 1.91]. This is distinct from the
primary paired \(t\) intervals.

\begin{table}[H]
\centering
\footnotesize
\caption{Archived architecture transfer on three-stage ArXiv. Each entry reports Macro-F1 / Accuracy averaged over nine memories and three student seeds per memory.}
\label{tab:sage}
\begin{tabular*}{\linewidth}{@{\extracolsep{\fill}}lrr@{}}
\toprule
Method & GCN & GraphSAGE\\
\midrule
CaT & 34.98 / 60.41 & 50.24 / 70.78\\
GECC & 52.68 / 70.94 & 52.60 / 70.62\\
\method{} & 52.91 / 71.15 & 53.93 / 71.27\\
\bottomrule
\end{tabular*}
\end{table}

\subsection{Support and Text Controls}

The comparison of shared and separate assignments matches initialization,
objectives, residuals, and construction budget. Sharing changes F1 by
+0.04 on three-stage ArXiv, -0.37 on SportsFit, and -1.92 on BookHis.
Their paired 95\% intervals are [-0.43, 0.50], [-1.21, 0.46], and
[-5.02, 1.17], respectively. One \(n_t\times m_t\) score matrix replaces two.
\autoref{tab:assignment_cost} and~\autoref{tab:arxiv_archived_controls}
report F1 and construction cost for these runs. The smaller parameterization
and measured cost reductions are distinct from the feature construction
effects in \autoref{tab:direct_features}.

\autoref{tab:support_diagnostics} describes the learned assignments and
stored text. Hard slot coverage is the fraction of slots selected by at
least one source node's largest row assignment. Normalized row entropy
divides each node's assignment entropy by the logarithm of its class's
slot count and averages over nodes, using zero for a class having one slot.
The unique text ratio is the number of distinct stored strings divided by
the number of slots.
Pairwise text cosine averages similarities between different selected text
embeddings among the first 512 slots, or all slots when fewer are available.
The soft/hard cosine gap is one minus the cosine similarity between each
slot's normalized soft text mixture and its selected source embedding.

\begin{table}[H]
\centering
\footnotesize
\caption{Support and text diagnostics for newly constructed blocks. Values average the recorded stages and seeds. ArXiv uses three stages.}
\label{tab:support_diagnostics}
\begin{tabular*}{\linewidth}{@{\extracolsep{\fill}}lrrr@{}}
\toprule
Diagnostic & ArXiv & SportsFit & BookHis\\
\midrule
Hard slot coverage & 1.0000 & 0.9998 & 1.0000\\
Normalized row entropy & 0.1200 & 0.2385 & 0.0827\\
Unique text ratio & 1.0000 & 0.9995 & 0.9990\\
Pairwise text cosine & 0.0947 & 0.3988 & 0.1194\\
Soft/hard cosine gap & 0.3853 & 0.2037 & 0.4693\\
\bottomrule
\end{tabular*}
\end{table}

Nearly all slots receive hard support and most stored texts are distinct.
The nonzero cosine gap shows that selecting a real text still differs from
retaining the full soft mixture.

\begin{table}[H]
\centering
\footnotesize
\caption{Text initialization and feedback controls. Macro-F1 on one order
and three construction seeds, each averaging three student model seeds.}
\label{tab:text_feedback}
\begin{tabular*}{\linewidth}{@{\extracolsep{\fill}}lrrr@{}}
\toprule
Constructor & ArXiv & SportsFit & BookHis\\
\midrule
\method{} & 52.96 & 80.02 & 48.53\\
No text initialization & 53.52 & 77.50 & 46.79\\
No pair feedback & 53.13 & 80.30 & 48.14\\
Both off & 53.32 & 77.64 & 47.87\\
\bottomrule
\end{tabular*}
\end{table}

Removing text initialization sets the text weight in prototypes and initial
scores to zero. Removing pair feedback sets the coefficient of the proxy's
graph and text pairing loss to zero, retaining prompt supervision.
The paired intervals for removing this feedback all cross zero.
The factorial does not show a general positive interaction.

\begin{table}[H]
\centering
\footnotesize
\caption{Text selection controls with fixed features. Entries report Macro-F1. Only stored text sources change, with a common limit of 320 characters.}
\label{tab:text_budget}
\begin{tabular*}{\linewidth}{@{\extracolsep{\fill}}lrrr@{}}
\toprule
Text selection & ArXiv & SportsFit & BookHis\\
\midrule
Primary & 52.96 & 80.09 & 47.70\\
Nearest-1 & 53.14 & 79.68 & 47.86\\
Nearest-5 & 52.59 & 80.75 & 47.98\\
Random-1 & 53.23 & 80.33 & 47.56\\
Random-5 & 53.01 & 80.65 & 46.56\\
\bottomrule
\end{tabular*}
\end{table}

Nearest-1 and Nearest-5 select one or five source texts by feature similarity
within the slot's class. Random-1 and Random-5 select the corresponding
number uniformly within that class. SportsFit and BookHis use reconstructed
features under the same protocol, compared with the primary text rule in
that reconstruction. The controls show that learned text selection does not
consistently improve F1. Earlier controls storing a soft embedding or one
shared text per class also vary by dataset. These interventions change the
stored text while keeping the learned features fixed.

\subsection{Resource Phases and Topology Scope}

\begin{table}[H]
\centering
\footnotesize
\setlength{\tabcolsep}{3pt}
\caption{Earlier resource audit with evaluation caches for the full graph resident.
Construction columns are peak GiB / seconds.
Student training and original-graph inference report peak MiB and have the
same allocations across matched methods.}
\label{tab:phase_resources}
\begin{tabular*}{\linewidth}{@{\extracolsep{\fill}}lrrr@{}}
\toprule
Phase / method & ArXiv & SportsFit & BookHis\\
\midrule
Construction: \method{} & 11.57/245.9 & 23.44/173.3 & 2.77/65.4\\
Construction: CaT & 0.85/78.2 & 1.05/173.1 & 0.45/28.3\\
Construction: GECC & 0.83/190.5 & 0.98/725.7 & 0.38/57.6\\
Student training (MiB) & 508.3 & 511.7 & 201.0\\
Inference (MiB) & 1678.9 & 1710.6 & 486.2\\
\bottomrule
\end{tabular*}
\end{table}

ArXiv uses final memories from the three-stage main schedule. SportsFit and BookHis use reconstructed
memories under the same protocol, whose hashes differ from the originals. These measurements diagnose cost
phases and do not overwrite primary performance. \autoref{tab:primary_time}
reports construction times averaged over nine combinations of class order and construction seed. They exclude loading and
common text caching and were collected with concurrent
jobs, so they are not isolated throughput benchmarks. Full-history training
resources are measured separately in \autoref{tab:student_resource_audit}.
The loader keeps features and normalized adjacency for the full graph on the GPU for evaluation. These are
outside the constructor's inputs but included in allocator peaks. Logical
access restricted to the current increment therefore does not imply that
only that increment resides in GPU memory.

\begin{table}[H]
\centering
\footnotesize
\setlength{\tabcolsep}{3pt}
\caption{Isolated initial-stage audit of the sparse edge-free operator.
The objective, Sinkhorn configuration, and proxy are unchanged. Only the
implementation of the self-loop matrix--vector product changes.}
\label{tab:sparse_edgefree_cost}
\begin{tabular*}{\linewidth}{@{\extracolsep{\fill}}lrrr@{}}
\toprule
Dataset & Dense GiB / s & Sparse GiB / s & Peak reduction\\
\midrule
ArXiv & 11.56 / 82.0 & 11.50 / 82.2 & 0.06 GiB\\
SportsFit & 23.44 / 123.1 & 23.33 / 123.1 & 0.11 GiB\\
BookHis & 1.17 / 20.0 & 1.17 / 19.6 & 0.001 GiB\\
\bottomrule
\end{tabular*}
\end{table}

The sparse operator passed isolated checks of forward values and first and
second derivatives, but its small peak reduction indicates
that dense support states and retained proxy intermediates dominate the
working set. We therefore use this option as an implementation optimization
and reproducibility check, not as evidence that the current constructor has
the lowest GPU memory cost.

\subsection{Optional Construction Storage}
\label{app:optional_storage}

The optional implementation stores only assignment entries between
current nodes and memory slots of the same class. Entries between
different classes are already constrained to zero. Storage by class
therefore reduces allocation without changing the objective, its
coefficients, or the dense denominator of the assignment anchor.
Evaluation caches are moved off the GPU during construction. The full
GM, DM, and propagation objectives and 30 outer updates are retained.
These implementation diagnostics use the development partition, one class
order, three construction seeds, and three student runs per memory.
ArXiv uses its three-stage class schedule, and SportsFit uses five stages.

\begin{table}[H]
\centering
\footnotesize
\setlength{\tabcolsep}{3pt}
\caption{Optional storage and cache offloading. Scores and peak allocated
GPU memory are from the same runs, averaged over three construction seeds.}
\label{tab:optional_storage}
\begin{tabular*}{\linewidth}{@{\extracolsep{\fill}}llrrr@{}}
\toprule
Dataset & Implementation & Macro-F1 & Accuracy & GiB\\
\midrule
ArXiv & Dense & 51.640 & 70.939 & 11.537\\
& Class blocks & 51.660 & 70.903 & 5.252\\
SportsFit & Dense & 79.681 & 85.464 & 23.802\\
& Class blocks & 79.413 & 84.433 & 17.514\\
\bottomrule
\end{tabular*}
\end{table}

The combined change lowers mean peak allocation by 54.48\% on ArXiv
and 26.42\% on SportsFit. Cache offloading alone saves about 0.3 GiB.
The class blocks account for most of the reduction. The paired F1
differences are +0.021 points on ArXiv (95\% interval [-0.557, 0.598])
and -0.268 on SportsFit ([-2.042, 1.506]). A SportsFit construction loses
1.09 points. Although the mathematical objective is preserved, numerical
changes can affect the learned memory and discrete text choices.

\begin{table}[H]
\centering
\footnotesize
\setlength{\tabcolsep}{3pt}
\caption{Incremental construction versus construction from all observed
nodes in one call, using the optional implementation. Peak allocated GPU
memory averages three construction seeds. These modes have different total update counts.}
\label{tab:construction_history}
\begin{tabular*}{\linewidth}{@{\extracolsep{\fill}}lrrr@{}}
\toprule
Dataset & Incremental GiB & Full GiB & Reduction\\
\midrule
ArXiv & 5.252 & 8.939 & 41.25\%\\
SportsFit & 17.514 & 18.851 & 7.09\%\\
\bottomrule
\end{tabular*}
\end{table}

The full-history reference uses all final legal construction nodes and
edges with the same final class quotas, in one 30-update call.
Incremental construction uses three or five such calls. The comparison
therefore measures peak working allocation under the two access modes,
rather than equal-compute performance. SportsFit's initial stage already
contains 90.62\% of the full legal assignment entries, including the
dominant class block. This limits the peak reduction achievable by its
class partition. The CPU source loader remains resident in both modes.

\begin{table}[H]
\centering
\footnotesize
\caption{Mean construction time in seconds over nine combinations of class order and construction seed.}
\label{tab:primary_time}
\begin{tabular*}{\linewidth}{@{\extracolsep{\fill}}lrrr@{}}
\toprule
Method & ArXiv & SportsFit & BookHis\\
\midrule
\method{} & 243.3 & 170.3 & 68.9\\
GECC & 189.5 & 735.1 & 52.0\\
CaT & 76.1 & 174.6 & 28.6\\
\bottomrule
\end{tabular*}
\end{table}

The separate topology extension raises the area under the receiver operating characteristic (ROC)
curve, reported as the area under the curve (AUC), for a link probe between memory slots on Cora from
0.8325 to 0.8591 but lowers Citeseer from 0.8479 to 0.8388.
The probe measures how well similarities between slots distinguish links.
It does not train a graph autoencoder or establish general link prediction performance.

\subsection{Construction Update Sensitivity}
\label{app:update_sensitivity}

\begin{figure*}[t]
\centering
\includegraphics[width=0.98\textwidth]{construction_loss_traces.png}
\caption{Construction trajectories without GM on a logarithmic loss scale,
separated by dataset and stage. Longer completed trajectories flatten, but
loss stabilization and downstream utility are different outcomes. Longer
SportsFit runs are unavailable because of OOM. ArXiv uses the three-stage class schedule.}
\Description{Nine small line plots show the construction objective against
updates for three datasets and three stages. Completed ArXiv and BookHis
runs flatten after the early decrease. This development diagnostic removes
gradient matching and does not establish convergence of the complete method.}
\label{fig:construction_loss}
\end{figure*}

The development diagnostic holds access, prefix, rate, encoder, seeds, and
student model budget fixed while varying outer updates. It removes GM because
the complete objective with second derivatives can exceed available GPU memory.
It therefore measures the behavior of the reduced objective.
\autoref{fig:construction_loss} shows the observed trajectories.

\begin{table}[H]
\centering
\footnotesize
\setlength{\tabcolsep}{2pt}
\caption{Update sensitivity without GM. Scores average completed construction
seeds and three student seeds per construction. Counts differ across budgets.
Paired changes use matching class orders and construction seeds. Failures are shown.}
\label{tab:construction_steps}
\begin{tabular*}{\linewidth}{@{\extracolsep{\fill}}llrrrr@{}}
\toprule
Dataset & Steps & Macro-F1 & Accuracy & Time (s) & GiB\\
\midrule
ArXiv & 30 & 51.702 & 70.957 & 225.6 & 11.09\\
& 60 & 51.809 & 71.503 & 1525.3 & 16.31\\
& 100 & -- & -- & Failed$^a$ & --\\
SportsFit & 30 & 79.980 & 85.969 & 153.5 & 23.77\\
& 60/100 & -- & -- & OOM & --\\
BookHis & 30 & 48.020 & 82.244 & 61.5 & 2.75\\
& 60 & 45.685 & 82.343 & 282.0 & 5.35\\
& 100 & 46.356 & 82.174 & 722.8 & 5.80\\
\bottomrule
\end{tabular*}
\par\smallskip
\raggedright $^a$Sinkhorn numerical failure. OOM denotes an out-of-memory failure.
\end{table}

Increasing updates from 30 to 60 changes F1 by -0.142 points on ArXiv and
-2.249 on BookHis for matched runs. Increasing from 30 to 100 changes
BookHis F1 by -1.578 points.
A practical plateau requires less than 1\% relative absolute objective
change in each of two consecutive 10-update windows. Twenty of 65 completed
stage traces satisfy it. In the 100-update BookHis runs, all five stages
satisfy it in both windows 80--90 and 90--100. This operational diagnostic
does not imply zero gradients or plateau at 30 updates on every stage.

